\chapter{Typed Composition and the Category Laws}
\label{ch04:main}

The previous chapter fixed the object layer: a named finite carrier with an
explicitly ordered element sequence, whose identity is the pair of name and
ordered carrier. This chapter builds the arrow layer on top of it and proves
the central structural fact of the kernel: the objects and arrows form a
category, with identities that are neutral and a composition that is
associative wherever it is defined. The qualification ``wherever it is
defined'' is not a formality. Composition in the kernel is a \emph{partial}
operation, deliberately so, and the precise condition under which it is
defined --- declared target equals declared source, as objects, meaning both
name and ordered carrier --- is the single most consequential design decision
in the entire system. Everything downstream, from the endomap/path separation
of Chapters~\ref{ch05:main} and~\ref{ch06:main} to the lift functor and the
normal-form machinery, inherits its soundness from the typing discipline
established here.

The chapter has four movements. First, we state the definitions of identity
and composition exactly as the kernel implements them, and we spell out what
the definedness condition does and does not permit
(Section~\ref{ch04:sec:defs}). Second, we prove the category-laws theorem in
full, pointwise, with no appeal to imported category-theoretic machinery
(Section~\ref{ch04:sec:laws}); the proof is elementary, and its elementarity
is part of the point, because an executable kernel should rest on arguments a
reviewer can check by hand. Third, we state the mismatch-rejection contract
--- the behavioral specification of what the implementation must do when the
definedness condition fails --- and we catalog the negative tests that
exercise it (Section~\ref{ch04:sec:contract}). Fourth, we give the
verification story in depth (Sections~\ref{ch04:sec:evidence}
and~\ref{ch04:sec:lean}): what the bounded exhaustive suites actually
enumerate, including the exact counting behind the $19{,}683$ associativity
triples and the $499$ typed maps; what bounded exhaustive testing does and
does not establish, situated against the small-scope hypothesis of
Jackson~\cite{jackson}; and what the pinned Lean formal seed kernel-checks,
including the axiom-audit discipline and the honest statement of the open
gap between the Python implementation and the Lean model.

Throughout, claim status is stated inline. The category laws carry a complete
informal proof (this chapter), bounded exhaustive executable evidence, and
exact Lean lemmas for the Lean model of the definitions. None of that is
proof-assistant verification of the running Python code; the equivalence of
the two artifacts is an open obligation, recorded as such in the ledger of
Appendix~B, and we will not blur that line.

\section{Identity and Composition}
\label{ch04:sec:defs}

\subsection{Arrows recalled}

We briefly recall the arrow definition from Chapter~\ref{ch03:main} in the
form we need here, since every proof in this chapter manipulates it directly.

\begin{definition}[Arrow judgment]
\label{ch04:def:arrow}
An \emph{arrow judgment} $f : X \to Y$ consists of three data: a declared
source object $X$, a declared target object $Y$, and a total map on standard
indices, represented as the tuple
\[
  \bigl(f(0), f(1), \dots, f(|X|-1)\bigr),
\]
where each $f(i)$ is a standard index of $Y$, that is, an integer with
$0 \le f(i) < |Y|$. The tuple has exactly $|X|$ entries: one image for each
source element, no more and no fewer. Totality is enforced at construction;
there is no notion of a partially defined arrow in the kernel.
\end{definition}

\begin{definition}[Extensional arrow equality]
\label{ch04:def:extensional}
Two arrows are \emph{extensionally equal} if and only if their declared
sources are equal as objects, their declared targets are equal as objects,
and their mapping tuples are equal as tuples. The presentation name attached
to an arrow for display and debugging purposes participates in neither
equality nor hashing.
\end{definition}

In the implementation (\texttt{category.py}), extensional equality is
realized by excluding the name field from the generated comparison of the
frozen dataclass \texttt{FiniteArrow}; the triple of source, target, and
mapping tuple is exposed as the arrow's \emph{denotation key}. We return to
the consequences of name-exclusion in Chapter~\ref{ch05:main}, where
presentation names for large powers are deliberately compressed; the
compression is sound precisely because names are extensionally inert.

\begin{notation}
\label{ch04:not:hom}
We write $\Hom(X, Y)$ for the set of arrow judgments from $X$ to $Y$. Since
an arrow is determined by its mapping tuple once $X$ and $Y$ are fixed, and
each of the $|X|$ tuple positions ranges over $|Y|$ values,
\[
  |\Hom(X, Y)| = |Y|^{|X|},
\]
with the convention $|Y|^{0} = 1$ (the unique empty map out of an empty
source) and $0^{|X|} = 0$ for $|X| \ge 1$ (no total map from a nonempty
source into an empty target). These conventions are not stylistic: the
bounded enumeration counts of Section~\ref{ch04:sec:evidence} depend on
them, and the executable suite includes the empty-carrier cases explicitly.
\end{notation}

\subsection{Identity}

\begin{definition}[Identity arrow]
\label{ch04:def:identity}
For each object $X$, the \emph{identity arrow} $\id_X : X \to X$ is the
arrow whose declared source and declared target are both $X$ and whose
mapping tuple is $(0, 1, \dots, |X|-1)$, that is, $\id_X(i) = i$ for every
standard index $i$.
\end{definition}

Note what the definition quantifies over: there is one identity per
\emph{object}, not one identity per carrier. Two objects with the same
ordered carrier but different names --- say
$X = (\texttt{register\_a}, (0, 1))$ and
$Y = (\texttt{register\_b}, (0, 1))$ --- have distinct identity arrows
$\id_X$ and $\id_Y$, because the arrows' declared sources differ as objects.
This is the arrow-level echo of name-inclusive object identity, and it is
what makes the definedness condition below strict enough to be useful.

\subsection{Composition and its definedness condition}

\begin{definition}[Composition]
\label{ch04:def:comp}
Let $f : X \to Y$ and $g : Y' \to Z$ be arrows. The composite $g \circ f$ is
\emph{defined if and only if $Y = Y'$ as objects} --- that is, the declared
target of $f$ and the declared source of $g$ agree in both name and ordered
carrier. When defined, $g \circ f : X \to Z$ is the arrow with declared
source $X$, declared target $Z$, and mapping tuple
\[
  (g \circ f)(i) = g\bigl(f(i)\bigr)
  \qquad \text{for every standard index } i \text{ of } X.
\]
When $Y \ne Y'$, the composite does not exist, and the implementation raises
a typed error (Section~\ref{ch04:sec:contract}). No coercion of any kind is
performed.
\end{definition}

The phrase ``no coercion of any kind'' compresses three distinct refusals,
each of which blocks a distinct failure mode. It is worth separating them,
because informal operator notations typically permit at least one of the
three, and each permission reintroduces a bug class the kernel exists to
exclude.

\begin{ruleenv}[No cardinality coercion]
\label{ch04:rule:card}
If $|Y| = |Y'|$ but $Y \ne Y'$, the composite $g \circ f$ is undefined. Equal
cardinality means only that a bijection \emph{exists} between the carriers;
it does not select one, and any silent selection (by position, by element
equality, or by hash order) would smuggle an unrecorded arrow into the
composite. If the caller intends a particular bijection, the caller must
construct it as an explicit arrow $b : Y \to Y'$ and write
$g \circ b \circ f$, at which point the bijection is part of the auditable
term.
\end{ruleenv}

\begin{ruleenv}[No carrier-only coercion]
\label{ch04:rule:carrier}
If $Y$ and $Y'$ have equal ordered carriers but different names, the
composite is undefined. Carrier equality without name equality means the two
objects are extensionally alike but were declared with distinct intent ---
two registers that both carry $(0, 1)$, two stages of a workflow that both
carry the same label set. Composing across them silently would erase the
declared distinction, and the entire point of name-inclusive identity
(Chapter~\ref{ch03:main}) is that declared distinctions are load-bearing. The
sanctioned bridge is again explicit: a caller who asserts that the two
objects should communicate constructs a named transport arrow between them
and includes it in the term.
\end{ruleenv}

\begin{ruleenv}[No name-only matching]
\label{ch04:rule:name}
Symmetrically, if $Y$ and $Y'$ share a name but differ in carrier --- an
error state that can arise when two modules independently declare an object
called \texttt{state} with different element orders --- the composite is
undefined. Object equality is the conjunction of name equality and carrier
equality; neither conjunct alone suffices. This refusal is what prevents a
renamed or reordered carrier from silently changing the meaning of every
downstream index.
\end{ruleenv}

\begin{remark}[Partiality is the feature]
\label{ch04:rem:partiality}
It is tempting to regard partial composition as an inconvenience to be
engineered away --- for instance, by embedding all carriers into one large
universal carrier and making every composite defined. The kernel refuses
this on principle. A total composition over a universal carrier is exactly
the untyped setting whose failure modes motivated the system: every
expression composes, including the meaningless ones, and mistypings surface
(if at all) as wrong answers rather than as errors. Partiality localizes the
failure to the construction site of the illegal term, with the two object
identities in hand, which is the earliest and cheapest place to fail.
\end{remark}

\begin{example}[A composite that exists]
\label{ch04:ex:good}
Let $X = (\texttt{stage\_in}, (a, b, c))$,
$Y = (\texttt{stage\_mid}, (u, v))$, and
$Z = (\texttt{stage\_out}, (p, q, r))$. Take $f : X \to Y$ with mapping
$(0, 0, 1)$ --- so $a \mapsto u$, $b \mapsto u$, $c \mapsto v$ --- and
$g : Y \to Z$ with mapping $(2, 0)$ --- so $u \mapsto r$, $v \mapsto p$.
The middle objects agree exactly, so $g \circ f : X \to Z$ is defined, and
its mapping tuple is
\[
  \bigl(g(f(0)), g(f(1)), g(f(2))\bigr) = \bigl(g(0), g(0), g(1)\bigr)
  = (2, 2, 0),
\]
that is, $a \mapsto r$, $b \mapsto r$, $c \mapsto p$.
\end{example}

\begin{example}[A composite that is refused]
\label{ch04:ex:bad}
With $X, Y, Z$ as above, let $Y'' = (\texttt{stage\_mid\_v2}, (u, v))$ ---
same ordered carrier as $Y$, different name --- and let $g'' : Y'' \to Z$
have the same mapping tuple $(2, 0)$ as $g$. The expression $g'' \circ f$ is
undefined: $f$ declares target $Y$, $g''$ declares source $Y''$, and
$Y \ne Y''$ because their names differ. The mapping tuples of $g$ and $g''$
are identical; the refusal is purely a matter of declared type. This is
Rule~\ref{ch04:rule:carrier} in action, and it is a genuine negative test in
the executable suite, not a hypothetical.
\end{example}

\subsection{Implementation fidelity: the tuple formula}
\label{ch04:sec:tuple}

Definition~\ref{ch04:def:comp} is phrased in terms of the abstract function
application $g(f(i))$. The implementation works with mapping tuples, and it
is worth writing down --- and proving --- the tuple-level formula it
computes, both because the executable tests compare tuples and because the
Lean model's composition is stated at exactly this level.

\begin{proposition}[Tuple form of composition]
\label{ch04:prop:tuple}
Let $f : X \to Y$ have mapping tuple $(a_0, \dots, a_{|X|-1})$ and let
$g : Y \to Z$ have mapping tuple $(b_0, \dots, b_{|Y|-1})$. Then the
composite $g \circ f$ of Definition~\ref{ch04:def:comp} has mapping tuple
\[
  \bigl(b_{a_0},\, b_{a_1},\, \dots,\, b_{a_{|X|-1}}\bigr),
\]
that is, the tuple obtained by indexing $g$'s tuple at each entry of $f$'s
tuple, in order. Every index operation involved is in range.
\end{proposition}

\begin{proof}
By Definition~\ref{ch04:def:arrow}, $f(i) = a_i$ and $g(j) = b_j$ for
standard indices $i$ of $X$ and $j$ of $Y$. Hence
$(g \circ f)(i) = g(f(i)) = g(a_i) = b_{a_i}$ for each $i$, which is the
displayed tuple read entrywise. Well-formedness of $f$ gives
$0 \le a_i < |Y|$, so $b_{a_i}$ is a legal lookup into a tuple of length
$|Y|$; well-formedness of $g$ gives $0 \le b_{a_i} < |Z|$, so every entry
of the result is a standard index of $Z$.
\end{proof}

The formula makes the computational cost visible --- composition is a
single pass over the source tuple, $\Theta(|X|)$ time with no search and
no allocation beyond the result --- and it makes the direction convention
unambiguous: the \emph{outer} arrow's tuple is indexed by the \emph{inner}
arrow's entries. Transposing the two tuples is the classic implementation
bug for this operation; it happens to produce a well-formed arrow whenever
$|X| = |Z|$ (so no validation catches it), and it is exactly the kind of
defect the exhaustive associativity sweep of
Section~\ref{ch04:sec:assoc-count} would expose, since the transposed
formula is not associative against the correct one. A worked instance of
the formula on Example~\ref{ch04:ex:good}: $f$'s tuple is $(0, 0, 1)$,
$g$'s tuple is $(2, 0)$, and indexing $g$'s tuple at $0, 0, 1$ yields
$(2, 2, 0)$, as computed there.

\begin{remark}[Where the definedness check sits]
\label{ch04:rem:checksite}
In the implementation, the object-equality comparison between the declared
target of the inner arrow and the declared source of the outer arrow is
the first statement of the composition function, executed before the tuple
of Proposition~\ref{ch04:prop:tuple} is built. The ordering matters for
the error contract of Section~\ref{ch04:sec:contract}: the check must not
be reachable only through some later validation failure (an out-of-range
index, a length mismatch) that would report the wrong cause. Because
carrier sizes can coincide while objects differ, a composition across
mismatched middles could produce an \emph{apparently} well-formed tuple;
only the up-front object comparison rejects it for the right reason.
\end{remark}

\section{The Category-Laws Theorem}
\label{ch04:sec:laws}

We now prove that the definitions of Section~\ref{ch04:sec:defs} satisfy the
category axioms. The statement is deliberately phrased in terms of the
concrete data --- declared sources, declared targets, mapping tuples ---
rather than by invoking the abstract definition and checking it off, because
the concrete phrasing is what the executable tests and the Lean lemmas
actually verify, clause by clause.

\begin{theorem}[Category laws]
\label{ch04:thm:laws}
Under Definitions~\ref{ch04:def:arrow}, \ref{ch04:def:identity},
and~\ref{ch04:def:comp}:
\begin{enumerate}
\item \textup{(Typing of identities)} For every object $X$, $\id_X$ is an
arrow from $X$ to $X$.
\item \textup{(Typing of composites)} Whenever $f : X \to Y$ and
$g : Y \to Z$, the composite $g \circ f$ is defined and is an arrow from $X$
to $Z$.
\item \textup{(Right neutrality)} For every $f : X \to Y$, the composite
$f \circ \id_X$ is defined and $f \circ \id_X = f$.
\item \textup{(Left neutrality)} For every $f : X \to Y$, the composite
$\id_Y \circ f$ is defined and $\id_Y \circ f = f$.
\item \textup{(Associativity)} Whenever $f : X \to Y$, $g : Y \to Z$, and
$h : Z \to W$, all four composites $g \circ f$, $h \circ g$,
$(h \circ g) \circ f$, and $h \circ (g \circ f)$ are defined, and
\[
  (h \circ g) \circ f = h \circ (g \circ f).
\]
\end{enumerate}
Consequently the finite objects and arrow judgments form a category, which
we denote $\cC$ in this chapter.

\textup{Status: complete informal proof below; bounded exhaustive executable
evidence (Section~\ref{ch04:sec:evidence}); exact Lean lemmas for the Lean
model (Section~\ref{ch04:sec:lean}); Python-to-Lean equivalence open.}
\end{theorem}

\begin{proof}
We verify each clause against the definitions, pointwise on standard
indices. Recall that by Definition~\ref{ch04:def:extensional}, two arrows
are equal exactly when their declared sources, declared targets, and mapping
tuples agree; so each equality claim below reduces to three componentwise
checks, of which the first two are bookkeeping and the third is the
pointwise computation.

\emph{Clause 1.} By Definition~\ref{ch04:def:identity}, $\id_X$ has declared
source $X$, declared target $X$, and mapping tuple $(0, \dots, |X|-1)$.
Every entry $i$ of the tuple satisfies $0 \le i < |X|$, so the tuple is a
valid total map into the target $X$, and the tuple length equals $|X|$, the
size of the source. Hence $\id_X$ is a well-formed arrow $X \to X$.

\emph{Clause 2.} Suppose $f : X \to Y$ and $g : Y \to Z$; the definedness
condition of Definition~\ref{ch04:def:comp} holds because the declared
target of $f$ and the declared source of $g$ are literally the same object
$Y$. We must check that the composite's mapping tuple is well-formed. Its
length is $|X|$, one entry per standard index of $X$. For each such index
$i$: $f(i)$ is a standard index of $Y$ because $f$ is a well-formed arrow
into $Y$, so $0 \le f(i) < |Y|$; therefore $f(i)$ is a legal argument
position for the mapping tuple of $g$, whose length is $|Y|$; and
$g(f(i))$ is then a standard index of $Z$ because $g$ is a well-formed arrow
into $Z$. Hence every entry of the composite tuple lies in
$\{0, \dots, |Z|-1\}$, and $g \circ f$ is a well-formed arrow $X \to Z$.
Note that this clause is where totality earns its keep: because $f$ is total
and lands inside $Y$'s index range, and $g$ is total on exactly that range,
the composite is total by construction, with no possibility of an
out-of-range lookup.

\emph{Clause 3.} The composite $f \circ \id_X$ is defined because the
declared target of $\id_X$ is $X$, which equals the declared source of $f$.
Its declared source is that of $\id_X$, namely $X$; its declared target is
that of $f$, namely $Y$; both agree with the corresponding data of $f$. For
the mapping tuple, fix any standard index $i$ of $X$:
\[
  (f \circ \id_X)(i) = f\bigl(\id_X(i)\bigr) = f(i),
\]
using $\id_X(i) = i$. Since $i$ was arbitrary, the mapping tuples agree
entrywise, hence as tuples. All three components of extensional equality
hold, so $f \circ \id_X = f$.

\emph{Clause 4.} The composite $\id_Y \circ f$ is defined because the
declared target of $f$ is $Y$, the declared source of $\id_Y$. Its declared
source is $X$ and its declared target is $Y$, agreeing with $f$. Pointwise,
for each standard index $i$ of $X$:
\[
  (\id_Y \circ f)(i) = \id_Y\bigl(f(i)\bigr) = f(i),
\]
using $\id_Y(j) = j$ at $j = f(i)$, which is legitimate because Clause~2's
argument shows $f(i)$ is a standard index of $Y$. Hence
$\id_Y \circ f = f$.

\emph{Clause 5.} First, definedness of all four composites. $g \circ f$ is
defined (target of $f$ is $Y$, source of $g$ is $Y$) and has type $X \to Z$
by Clause~2. $h \circ g$ is defined (target of $g$ is $Z$, source of $h$ is
$Z$) and has type $Y \to W$. Then $(h \circ g) \circ f$ is defined because
the target of $f$ is $Y$ and the source of $h \circ g$ is $Y$; and
$h \circ (g \circ f)$ is defined because the target of $g \circ f$ is $Z$
and the source of $h$ is $Z$. Both outer composites have declared source $X$
and declared target $W$, so the first two components of extensional
equality agree.

For the mapping tuples, fix a standard index $i$ of $X$. Unfolding
Definition~\ref{ch04:def:comp} twice on each side:
\[
  \bigl((h \circ g) \circ f\bigr)(i)
  = (h \circ g)\bigl(f(i)\bigr)
  = h\Bigl(g\bigl(f(i)\bigr)\Bigr),
\]
\[
  \bigl(h \circ (g \circ f)\bigr)(i)
  = h\Bigl((g \circ f)(i)\Bigr)
  = h\Bigl(g\bigl(f(i)\bigr)\Bigr).
\]
The two right-hand sides are the same expression --- three nested total
lookups $h(g(f(i)))$ --- so the tuples agree at every index, hence as
tuples. Therefore $(h \circ g) \circ f = h \circ (g \circ f)$.

All five clauses hold, and the data $(\text{objects}, \Hom, \id, \circ)$
satisfy the axioms of a category in the standard sense
\cite{maclanecwm}.
\end{proof}

\begin{remark}[Why the proof is this short]
\label{ch04:rem:short}
The proof is five paragraphs of index-chasing, and a reader familiar with
$\Set$ will recognize it as the standard argument that finite sets and total
functions form a category, decorated with bookkeeping about declared names.
That is exactly right, and it is worth being explicit about what the
decoration buys. The mathematical content --- functions compose
associatively, identities are neutral --- is classical and utterly standard
\cite{maclanecwm}. The kernel's contribution is not the theorem but the
\emph{definedness discipline} under which the theorem is stated: composition
is partial, the partiality tracks declared intent (names) and not merely
extension (carriers), and the associativity clause is proved only for
triples that are composable under that strict discipline. The theorem's
hypotheses are the contract; the proof is routine once the contract holds.
The system's engineering claim is that the contract is enforced at
construction time, every time, which is a statement about the
implementation, verified by the evidence of
Sections~\ref{ch04:sec:evidence} and~\ref{ch04:sec:lean} to the bounded
extents stated there.
\end{remark}

\begin{corollary}[Uniqueness of identities]
\label{ch04:cor:unique-id}
For each object $X$, $\id_X$ is the unique arrow $e : X \to X$ such that
$f \circ e = f$ for every $f$ with source $X$ and $e \circ g = g$ for every
$g$ with target $X$.
\end{corollary}

\begin{proof}
Existence is Theorem~\ref{ch04:thm:laws}, Clauses~3 and~4. For uniqueness,
suppose $e : X \to X$ satisfies the right-neutrality property. Instantiate
it at $f = \id_X$: then $\id_X \circ e = \id_X$. But by left neutrality of
$\id_X$ (Clause~4 applied to $e$), $\id_X \circ e = e$. Hence
$e = \id_X$.
\end{proof}

\begin{corollary}[Endomorphism monoids]
\label{ch04:cor:monoid}
For each object $X$, the set $\Hom(X, X)$ of endo-arrows on $X$ is a monoid
under composition, with identity element $\id_X$, of cardinality
$|X|^{|X|}$.
\end{corollary}

\begin{proof}
Composition of two endo-arrows on $X$ is always defined (the middle object
is $X$ on both sides) and yields an endo-arrow on $X$
(Theorem~\ref{ch04:thm:laws}, Clause~2), so $\Hom(X, X)$ is closed under
$\circ$. Associativity is Clause~5 restricted to endo-arrows; neutrality of
$\id_X$ is Clauses~3 and~4. The cardinality is
Notation~\ref{ch04:not:hom} with $Y = X$.
\end{proof}

Corollary~\ref{ch04:cor:monoid} is the doorway to
Chapter~\ref{ch05:main}: the monoid structure on $\Hom(X, X)$ is what makes
powers of a single endomap meaningful, and its \emph{absence} for arrows
whose source and target differ is what makes powers of heterogeneous data
meaningless. The kernel promotes this mathematical asymmetry into a type
asymmetry, and Chapters~\ref{ch05:main} and~\ref{ch06:main} develop the two
sides.

\begin{remark}[What the category is not]
\label{ch04:rem:notset}
$\cC$ is not $\Set$, and not even the category of finite sets, in one
substantive respect: its objects are \emph{named} carriers, so the same
finite carrier appears as the underlying data of many non-isomorphic-by-fiat
objects --- or more precisely, many distinct objects that happen to be
isomorphic via the evident positional bijection but are not \emph{equal},
and between which no arrow exists until one is explicitly constructed. The
skeleton of $\cC$ is equivalent to the category of finite sets in the usual
way, and nothing in this monograph depends on any exotic property; the
naming layer is an engineering surface, not new mathematics. We flag this to
preempt a possible over-reading: no claim is made that name-inclusive
identity changes the categorical content. It changes which composites are
\emph{writable}, which is a statement about the term language, not about the
equivalence class of the category.
\end{remark}

\section{The Mismatch-Rejection Contract}
\label{ch04:sec:contract}

Theorem~\ref{ch04:thm:laws} governs the composites that exist. An executable
kernel must be equally precise about the composites that do not, because in
practice the value of a type discipline is measured at its boundary: what
happens when a caller writes the illegal term. This section states the
contract and catalogs the negative tests.

\subsection{The contract}

\begin{invariant}[Mismatch rejection]
\label{ch04:inv:mismatch}
For all arrows $f$ and $g$, the implementation of $g \circ f$ either
\begin{enumerate}
\item returns the arrow of Definition~\ref{ch04:def:comp}, in the case that
the declared target of $f$ equals the declared source of $g$ as objects
(name and carrier both); or
\item raises the dedicated typed error \texttt{CompositionTypeError},
carrying a message that names both offending arrows and their declared
types, in every other case.
\end{enumerate}
There is no third behavior: no silent coercion, no warning-and-proceed, no
best-effort composite over a common sub-carrier, no reordering of either
carrier to force a match.
\end{invariant}

Three properties of the contract deserve comment.

\emph{The error is typed.} \texttt{CompositionTypeError} is a distinct
exception class (a subclass of the host language's \texttt{TypeError}), not
a generic failure. Callers and test harnesses can therefore assert the
\emph{reason} for a failure, not merely its occurrence. A negative test that
catches a bare exception would pass even if the implementation failed for
the wrong reason --- an off-by-one in tuple validation, say --- and such a
test would be evidence of nothing. Every negative fixture in the suite
asserts the specific exception class.

\emph{The error is early.} The check runs before any part of the composite
tuple is constructed. There is no partially built arrow to leak, and the
failure site is the construction expression itself, with both declared types
in the error message. In a system whose purpose is auditability, the
distance between the mistake and its detection is a first-class quality
metric, and here that distance is zero.

\emph{The error is symmetric in the three refusals.}
Rules~\ref{ch04:rule:card}, \ref{ch04:rule:carrier},
and~\ref{ch04:rule:name} all funnel into the same check --- object equality
of the middle objects --- because object equality is the conjunction of name
and carrier equality. The implementation does not special-case ``close''
mismatches; a near-miss (equal carriers, different names) and a gross
mismatch (disjoint carriers of different sizes) are rejected by the same
comparison at the same site. This uniformity is deliberate: any
special-casing would be a coercion pathway waiting to be widened.

\subsection{The negative-test catalog}

The executable suite exercises the contract with a family of negative
fixtures, each targeting one way the definedness condition can fail. We
catalog them as checkclasses; Chapter~\ref{ch10:main} gives the mapping to
concrete test names.

\begin{checkclass}[Carrier-equal, name-distinct middle]
\label{ch04:check:nameonly}
Arrows $f : X \to Y$ and $g : Y'' \to Z$ with $Y$ and $Y''$ carrying equal
ordered carriers under different names, as in Example~\ref{ch04:ex:bad}.
Expected: \texttt{CompositionTypeError}. This is the sharpest fixture in the
family, because every coercion-friendly design passes it wrongly: any
implementation that compares carriers, cardinalities, or element sets ---
anything short of full object identity --- will accept the composite. The
fixture therefore discriminates the kernel's discipline from its most
plausible near-neighbors.
\end{checkclass}

\begin{checkclass}[Cardinality-equal, carrier-distinct middle]
\label{ch04:check:cardonly}
Middle objects of equal size with different element sequences (different
elements, or the same elements in a different order). Expected:
\texttt{CompositionTypeError}. The reordered-carrier subcase is the
determinism guard: two declarations of ``the same'' set in different orders
are different objects, and composing across them would silently permute
every downstream index.
\end{checkclass}

\begin{checkclass}[Gross type mismatch]
\label{ch04:check:gross}
Middle objects differing in both name and size. Expected:
\texttt{CompositionTypeError}. This is the base case, included so that the
family covers the full lattice of agreement patterns between name, carrier,
and cardinality.
\end{checkclass}

\begin{checkclass}[Malformed arrow construction]
\label{ch04:check:malformed}
Not a composition test but a prerequisite for the others: attempts to
construct arrows whose mapping tuple has the wrong length, contains an
out-of-range index, contains a non-integer (including the host language's
boolean values, which the implementation rejects explicitly even though the
host treats them as integers), or is supplied as an unordered container.
Expected: the construction-time \texttt{ValidationError}. These fixtures
ensure that every arrow reaching a composition site is well-formed, which is
the standing hypothesis of Theorem~\ref{ch04:thm:laws}.
\end{checkclass}

\begin{checkclass}[Path-adjacency mismatch]
\label{ch04:check:path}
The path constructor of Chapter~\ref{ch06:main} performs the same middle-object
check between each adjacent pair of a sequence of arrows; fixtures supply
sequences with one ill-typed adjacency and assert
\texttt{CompositionTypeError} naming the offending segment. We list it here
because it is the same contract, enforced at a second surface.
\end{checkclass}

\subsection{The agreement lattice}
\label{ch04:sec:lattice}

The negative fixtures are not an ad hoc list; they cover a small, complete
lattice of ways the middle objects can agree or disagree, and completeness
of the coverage is checkable by inspection of that lattice. The middle
objects $Y$ (declared target of the inner arrow) and $Y'$ (declared source
of the outer arrow) can agree or disagree in name and in ordered carrier;
carrier agreement further refines by whether the carriers at least have
equal cardinality. Since carrier equality implies cardinality equality,
the meaningful combinations are five:

\begin{center}
\small
\begin{tabular}{@{}lllll@{}}
\toprule
Names & Carriers & Cardinalities & Composite & Exercised by \\
\midrule
equal & equal & equal & \emph{defined} & positive suites
(Sections~\ref{ch04:sec:assoc-count}, \ref{ch04:sec:typed-count}) \\
distinct & equal & equal & refused & Checkclass~\ref{ch04:check:nameonly} \\
equal & distinct & equal & refused & Checkclass~\ref{ch04:check:cardonly} \\
distinct & distinct & equal & refused & Checkclass~\ref{ch04:check:cardonly} \\
any & distinct & distinct & refused & Checkclass~\ref{ch04:check:gross} \\
\bottomrule
\end{tabular}
\end{center}

Only the top row composes; the implementation reaches all five rows
through the single object-equality comparison, and the fixture families
place at least one test in each refused row. The third row --- equal
names, distinct carriers --- deserves the extra emphasis given to it by
Rule~\ref{ch04:rule:name}: it models the situation where two modules
independently declare an object under the same name with different element
orders, and it is the row where a name-keyed registry design (looking
objects up by name and trusting the hit) would silently compose across
genuinely different types. The kernel has no registry; objects are values,
and equality is structural over the full pair.

\begin{boundary}
\label{ch04:bound:negative}
The negative-test catalog establishes that the implementation rejects the
enumerated fixture families. It does not establish that the implementation
rejects \emph{every} ill-typed composite expressible in the host language:
that universal statement is a property of the code paths, argued informally
by inspection (there is exactly one composition function, and its first
action is the middle-object equality check) and not verified by proof
assistant. The distinction matters and is recorded in the ledger
(PO-TKS8-CORE-02, status \textsc{executable validation}).
\end{boundary}

\section{The Verification Story I: Bounded Exhaustive Testing}
\label{ch04:sec:evidence}

The kernel's executable evidence for Theorem~\ref{ch04:thm:laws} consists of
bounded exhaustive enumeration: not sampled property tests but complete
sweeps of every instance within a stated finite scope. This section explains
exactly what is enumerated, reproduces the counting arithmetic behind the
recorded case counts, and then confronts the epistemological question of
what such testing establishes.

\subsection{The associativity sweep: counting the $19{,}683$ triples}
\label{ch04:sec:assoc-count}

Fix one three-element object $X_3$ --- a single named carrier of size three.
By Corollary~\ref{ch04:cor:monoid}, the endo-arrows on $X_3$ form the full
transformation monoid on three points, of cardinality
\[
  |{\Hom(X_3, X_3)}| = 3^{3} = 27:
\]
each of the three source indices independently maps to one of three target
indices. Because every arrow in this family has source and target $X_3$,
\emph{every ordered triple} $(f, g, h)$ of them is composable in both
associations --- the definedness condition holds automatically at both
middle objects --- and so the associativity clause can be tested on all
triples without any composability filtering. The number of ordered triples
is
\[
  27^{3} = \bigl(3^{3}\bigr)^{3} = 3^{9} = 19{,}683.
\]
The suite constructs, for each of the $19{,}683$ triples, both
$(h \circ g) \circ f$ and $h \circ (g \circ f)$, and asserts their
extensional equality --- source, target, and full mapping tuple. This is the
recorded figure of $19{,}683$ associativity triples: it is the complete
associativity table of the transformation monoid $T_3$, checked in extension.

Two remarks on the choice of scope. First, why endo-arrows on one carrier:
because that is the unique way to make every triple composable, which turns
the sweep into a clean product count with no reachability analysis; the
typed (multi-carrier) sweep below covers the heterogeneous cases. Second,
why size three: size two gives only $4^{3} = 64$ triples, which is thin, and
size four gives $256^{3} = 16{,}777{,}216$ triples, which is affordable but
was not the recorded snapshot's scope; the recorded suite pinned
$19{,}683$. The count is exact and the enumeration is complete within its
scope; the scope is what it is, and Section~\ref{ch04:sec:smallscope}
addresses what that means.

\subsection{The typed sweep: counting the $499$ maps}
\label{ch04:sec:typed-count}

The second family exercises the typed discipline across \emph{distinct}
carriers: identity neutrality, typed composition, and mismatch rejection for
maps whose source and target sizes range through four, including the empty
carrier. The recorded count is $499$ typed maps, and the arithmetic behind
it is the complete enumeration of Notation~\ref{ch04:not:hom} over one
representative source-and-target pair per size combination:
\[
  \sum_{m=0}^{4} \sum_{n=0}^{4} n^{m} = 499,
\]
with the conventions $n^{0} = 1$ (including $0^{0} = 1$: the empty map from
the empty object to the empty object) and $0^{m} = 0$ for $m \ge 1$.
The row-by-row expansion makes the conventions visible:
\begin{center}
\begin{tabular}{@{}lrrrrrr@{}}
\toprule
Source size $m$ & $n{=}0$ & $n{=}1$ & $n{=}2$ & $n{=}3$ & $n{=}4$ & Row total \\
\midrule
$m = 0$ & $1$ & $1$ & $1$ & $1$ & $1$ & $5$ \\
$m = 1$ & $0$ & $1$ & $2$ & $3$ & $4$ & $10$ \\
$m = 2$ & $0$ & $1$ & $4$ & $9$ & $16$ & $30$ \\
$m = 3$ & $0$ & $1$ & $8$ & $27$ & $64$ & $100$ \\
$m = 4$ & $0$ & $1$ & $16$ & $81$ & $256$ & $354$ \\
\midrule
Total & & & & & & $499$ \\
\bottomrule
\end{tabular}
\end{center}
The presence of the $m = 0$ row and the $n = 0$ column is not pedantry.
Empty carriers are where off-by-one and vacuous-quantification bugs live:
the unique arrow out of the empty object composes with everything
typed-compatibly, its mapping tuple is empty, and an implementation that
mishandles zero-length tuples would fail exactly here and nowhere else. The
suite's inclusion of the degenerate sizes is part of its value, and the lift
functor's own sweep (Chapter~\ref{ch07:main}) makes the same choice for the
same reason.

For every enumerated map the suite checks the identity laws
($f \circ \id = f$ and $\id \circ f = f$, with the appropriate identities),
and for composable pairs among the enumerated maps it checks the typed
composition of Definition~\ref{ch04:def:comp}; the negative fixtures of
Section~\ref{ch04:sec:contract} run against representatives of the
mismatched combinations.

\subsection{Design alternative: sampled property testing}
\label{ch04:sec:sampled}

A reviewer familiar with property-based testing frameworks will ask why
the suite enumerates rather than samples: the standard tooling generates
random arrows, checks the law, shrinks counterexamples, and scales to
carriers of size fifty as easily as size four. The kernel's choice of
bounded exhaustion over sampling is deliberate, and the reasoning is
worth recording because it recurs in every later chapter's evidence
design.

First, coverage claims compose better than confidence claims. ``Every
triple through the stated bound passes'' is a statement another artifact
can cite and build on --- the Lean equivalence obligation, for instance,
could discharge part of its burden by replaying exactly the enumerated
family against a verified oracle. ``No counterexample found in $10^{5}$
random triples'' supports no such composition; its evidentiary content
depends on the generator's distribution, which is itself unverified code,
and distributional subtleties in random generation of structured data are
a classic source of vacuous test suites (generators that never produce
empty carriers, or never produce equal-carrier-distinct-name pairs,
silently test nothing about those cases).

Second, the failure modes plausible for this code are small-instance
failure modes. The composition path branches on nothing but tuple
lengths and index values; a defect that manifests only above size four
would need to be size-dependent in a way the code has no means to
express. This is the small-scope hypothesis \cite{jackson} applied with
knowledge of the implementation, which is the strongest form in which it
can be applied.

Third, determinism of the evidence artifact. The exhaustive suite is a
pure function of the code under test: two runs, on two machines, in two
decades, enumerate the identical case list in the identical order. A
sampled suite is a function of the code and a seed, and its transcript is
evidence only relative to that seed's provenance. In a system whose
entire object-layer design exists to eliminate interpreter-state
nondeterminism (Chapter~\ref{ch03:main}), reintroducing randomness in the
verification layer would be a discordant note.

None of this is an argument against sampling in general --- for spaces
that explode combinatorially, later chapters accept different trade-offs
--- but for the category laws the exhaustive design dominates: the space
is small, the enumeration is cheap, and the resulting claim is exact.

\subsection{What bounded exhaustive testing establishes}
\label{ch04:sec:smallscope}

We now state, carefully, the epistemic status of the numbers just derived.

Bounded exhaustive testing occupies a specific point between random testing
and proof. Unlike sampled property-based testing, it leaves no gaps
\emph{within its scope}: every instance up to the bound is checked, so any
counterexample of size at most the bound is found with certainty, not with
probability. Unlike proof, it says nothing beyond the bound: a defect that
first manifests on carriers of size five, or on a thirty-first arrow, or on
a triple whose middle carrier is empty in a configuration the sweep did not
include, is invisible to it.

The methodological warrant for taking small scopes seriously is the
\emph{small-scope hypothesis}, articulated by Jackson in the context of the
Alloy analyzer \cite{jackson}: a large fraction of the defects in structural
specifications can be exhibited by counterexamples of small size, so
exhaustive checking within a small bound catches a disproportionate share of
real bugs. For the properties at issue here the hypothesis is unusually
plausible, because the properties are \emph{uniform in the carrier}: the
associativity computation $h(g(f(i)))$ does not branch on carrier size, and
an implementation error in index composition --- composing tuples in the
wrong order, off-by-one in an index bound, mishandling the empty tuple ---
would almost certainly manifest already at sizes zero through four. The
sweeps are designed around exactly this: they saturate the small sizes
completely rather than sampling large ones sparsely.

But plausibility is not proof, and the hypothesis is a hypothesis. The
honest statement is the following, and the monograph will repeat it wherever
the numbers are cited:

\begin{boundary}
\label{ch04:bound:testing}
The $19{,}683$-triple and $499$-map sweeps establish that the implementation
satisfies the category laws and the typed-composition contract \emph{on
every instance within their stated bounds}, and nothing more. They are not
proof-assistant verification of the implementation; they are not evidence
about carriers of size five or beyond except via the small-scope heuristic,
which is an empirical regularity about bug distributions
\cite{jackson}, not a theorem; and passing them closes no open obligation in
the ledger. The mathematical theorem (Theorem~\ref{ch04:thm:laws}) is proved
for all sizes by the pointwise argument; the \emph{testing} speaks only to
whether the code implements the definitions, and only within scope.
\end{boundary}

It is worth dwelling on the division of labor in that boundary statement,
because it is easy to garble in summary. There are three artifacts: the
mathematical proof, the implementation, and the test suite. The proof is
complete and covers all finite sizes, but it is about the
\emph{definitions}, not the code. The tests connect the code to the
definitions, but only within bounds. The missing third leg --- a machine-
checked connection between code and definitions at all sizes --- is exactly
what the Lean seed begins and does not finish, and it is to the seed that we
now turn.

\section{The Verification Story II: The Lean Formal Seed}
\label{ch04:sec:lean}

\subsection{What kernel-checking means}

The kernel's formal artifact is a pinned Lean~4.32.1 development
\cite{lean4} containing, among other modules, \texttt{TypedCategory}: a Lean
model of the definitions of Section~\ref{ch04:sec:defs} together with
machine-checked proofs of the identity and associativity lemmas for that
model.

``Machine-checked'' has a precise meaning that is stronger than ``a tool
accepted the file,'' and the distinction is the entire value proposition.
Lean is architected around a small trusted kernel: the elaborator, the
tactic framework, and the standard automation are large and unverified, but
every proof they produce is re-checked, term by term, by the kernel, which
implements only the type theory's inference rules. A theorem ``kernel-checks''
when the fully elaborated proof term passes this final validation. Trust in
a kernel-checked theorem therefore reduces to trust in (i) the kernel
implementation, a deliberately small and heavily scrutinized artifact,
(ii) the consistency of the axioms the proof depends on, and (iii) the
fidelity of the \emph{statement} --- whether the formalized proposition says
what the informal one says. Points (i) and (ii) are standard and shared with
every Lean development; point (iii) is where all the residual risk
concentrates, and Section~\ref{ch04:sec:gap} takes it up.

\subsection{The axiom audit}

Lean proofs can depend on axioms: the standard classical trio
(\texttt{propext}, \texttt{Classical.choice}, \texttt{Quot.sound}), or worse,
project-declared axioms, which are simply unproved assumptions, and
\texttt{sorry}, which is an explicit hole. A development that quietly
depends on a project axiom proves nothing about anything; a development
audit must therefore inspect the transitive axiom dependencies of every
theorem it advertises.

The verifier script for the formal seed
(\texttt{scripts/verify\_psr8\_lean\_gate1.py}) mechanizes this audit. It
pins the toolchain (Lean~4.32.1, exact commit), checks source hashes so the
audited files are the shipped files, builds the development through Lake,
and then extracts the exact axiom-dependency report of each exported
theorem. The policy it enforces is maximally strict for this development:
the \emph{sole} permitted axiom is \texttt{propext} --- propositional
extensionality, the principle that logically equivalent propositions are
equal, one of Lean's standard axioms and among the least controversial ---
and the verifier \emph{rejects} any occurrence of a project-declared axiom
and any occurrence of \texttt{sorry}. The \texttt{TypedCategory} identity
and associativity lemmas pass under this policy: their proofs are
constructive index-chasing of exactly the shape of the proof of
Theorem~\ref{ch04:thm:laws}, and need nothing beyond \texttt{propext} (used,
where at all, in the ambient library plumbing rather than the combinatorial
core).

\begin{finding}
\label{ch04:find:lean}
The Lean model of Definitions~\ref{ch04:def:arrow}--\ref{ch04:def:comp}
satisfies the identity and associativity laws, as kernel-checked theorems
depending on no axiom beyond \texttt{propext}, with project axioms and
\texttt{sorry} mechanically excluded by the audit script, under the pinned
Lean~4.32.1 toolchain. \textup{Status: verified within the stated scope;
this finding is about the Lean model, not about the Python implementation.}
\end{finding}

\subsection{The open gap: Python-to-Lean equivalence}
\label{ch04:sec:gap}

The final clause of Finding~\ref{ch04:find:lean} is the load-bearing caveat
of this chapter, and we state it as its own obligation rather than a
footnote.

The Lean development formalizes a \emph{model} of the kernel's definitions:
Lean types for named finite objects and total index maps, a Lean composition
operation with the same definedness condition, Lean statements of the laws.
The Python implementation realizes the \emph{same definitions by intent} ---
the informal specification and this monograph's prose are the bridge --- but
no formal artifact connects the two. Nothing machine-checked establishes
that \texttt{compose} in \texttt{category.py} and the composition of the
Lean model compute the same function of their arguments, or that the
validation predicates coincide, or even that the two data representations
are in the evident bijection. The equivalence claim currently rests on
side-by-side inspection of two short pieces of code by human readers.

This is a familiar situation in verification practice --- it is the standard
``verified model versus running artifact'' gap --- and there are standard
ways to close it, each with costs: re-implement the kernel in Lean and
extract or transpile executable code from the proved artifact; or state the
Python semantics formally and prove refinement; or generate exhaustive
conformance vectors from the Lean model and replay them against Python,
which would still be testing, but testing against a verified oracle. None of
these has been done. The ledger records the gap as open
(PO-TKS8-CORE-01 carries \textsc{partial formal verification present}
precisely because the Lean side alone is verified; PO-TKS8-CORE-07 records
the formalization seed as partial and present), and no artifact of this
monograph should be summarized as ``the kernel is formally verified.'' The
kernel's category layer has: a complete informal proof about its
definitions; bounded exhaustive tests of its implementation; and a
kernel-checked formalization of its model under a one-axiom policy. Those
are three genuinely strong legs. The fourth --- proved equivalence of
implementation and model --- is open, and stating so precisely is a result
of this chapter, not a deficiency in it.

\section{Exercises}
\label{ch04:sec:exercises}

The following exercises are checkable entirely within the material of this
chapter; the first two are computations, the rest are short proofs or
design analyses. They exist because the chapter's claims are of the kind
best internalized by attempting to break them.

\begin{exercise}
\label{ch04:ex:compute}
Let $X$ carry four elements and let $f, g : X \to X$ have mapping tuples
$(1, 2, 3, 0)$ and $(0, 0, 2, 2)$ respectively. Compute the tuples of
$g \circ f$ and $f \circ g$ via Proposition~\ref{ch04:prop:tuple}, and
confirm they differ --- endomorphism monoids are not commutative in
general, which is why Corollary~\ref{ch05:cor:commute} (powers of a
\emph{single} arrow commute) has content.
\end{exercise}

\begin{exercise}
\label{ch04:ex:count5}
Redo the counting of Section~\ref{ch04:sec:typed-count} with the bound
raised from four to five and verify the total is
$\sum_{m,n \le 5} n^{m} = 4{,}425$. (This is the size the typed sweep
would have at scope five; the recorded suite pins scope four.)
\end{exercise}

\begin{exercise}
\label{ch04:ex:transposed}
Define the \emph{transposed} composition of tuples by indexing the inner
tuple at the outer tuple's entries (the bug of
Section~\ref{ch04:sec:tuple}). Exhibit three endo-arrows on a two-element
carrier for which the transposed operation fails associativity against
extensional equality, confirming that the $19{,}683$-triple sweep would
catch the bug at size three or below.
\end{exercise}

\begin{exercise}
\label{ch04:ex:uniqueid}
Show that Corollary~\ref{ch04:cor:unique-id} still holds if the
right-neutrality hypothesis is weakened to hold for the single arrow
$f = \id_X$. Conclude that the identity is determined by its behavior
against itself and one side of neutrality.
\end{exercise}

\begin{exercise}
\label{ch04:ex:coercion}
Suppose a hypothetical variant of the kernel added carrier-only coercion
(Rule~\ref{ch04:rule:carrier} dropped): composition is accepted whenever
the middle \emph{carriers} agree, names ignored. Show first that the
variant still satisfies all five clauses of
Theorem~\ref{ch04:thm:laws} --- tuple composition is associative and
unital regardless of names --- so no law-level test, exhaustive or not,
can distinguish the variant from the kernel. Then explain why the variant
nevertheless fails Checkclass~\ref{ch04:check:nameonly}, and conclude that
the negative fixtures, not the law sweeps, are the artifacts that pin the
kernel's typing discipline. (This is the precise sense in which the
mismatch-rejection contract is an independent specification rather than a
consequence of the category laws.)
\end{exercise}

\section{Summary of Claims and Their Status}
\label{ch04:sec:summary}

\begin{center}
\small
\begin{tabular}{@{}p{5.2cm}p{7.6cm}@{}}
\toprule
Claim & Status \\
\midrule
Category laws for the definitions (Theorem~\ref{ch04:thm:laws}) &
Complete informal pointwise proof, this chapter; standard mathematics
\cite{maclanecwm}. \\
Category laws in the implementation &
Bounded exhaustive evidence: $19{,}683$ associativity triples
($27^{3}$ over the $3^{3}$ endo-arrows on a three-carrier), $499$ typed maps
($\sum_{m,n \le 4} n^{m}$, empty carriers included); not verification
\cite{jackson}. \\
Mismatch rejection (Invariant~\ref{ch04:inv:mismatch}) &
Executable validation: typed error class, negative fixture families of
Section~\ref{ch04:sec:contract}; universal rejection argued by code
inspection only. \\
Category laws in the Lean model &
Kernel-checked under pinned Lean~4.32.1, \texttt{propext}-only axiom audit,
\texttt{sorry} rejected \cite{lean4}. \\
Python-to-Lean equivalence &
\textsc{Open}. No machine-checked connection between implementation and
model; recorded in the ledger. \\
\bottomrule
\end{tabular}
\end{center}

The next two chapters put the category to work. Chapter~\ref{ch05:main}
develops the one place where iteration is licensed --- the endomorphism
monoid of Corollary~\ref{ch04:cor:monoid} --- through to the finite
idempotence theorem and its verification. Chapter~\ref{ch06:main} develops
the heterogeneous side --- semigroupoid paths --- and analyzes, in three
extended case studies, the invalid inference that the separation between the
two is engineered to make unwritable.
